\documentclass[11pt]{article}

% ---------------------------------------------------------------------------
% Executable Conformance for a Secure-Multiparty Application
% FIRST DRAFT. Baseline artifact: commit 305a3a8 of smc-belief-tracking.
% Every substantive claim traces to a repository file (see \repo) or is
% flagged \open{...}. No research gap is silently filled.
% Venue class is [OPEN]; article is used so the draft compiles anywhere.
% ---------------------------------------------------------------------------

\usepackage[margin=1in]{geometry}
\usepackage{amsmath,amssymb,amsthm}
\usepackage{booktabs}
\usepackage{array}
\usepackage{tabularx}
\usepackage{xcolor}
\usepackage{listings}
\usepackage{enumitem}
\usepackage[hidelinks]{hyperref}
\usepackage{microtype}

% ---- helper commands -------------------------------------------------------
\newcommand{\code}[1]{\texttt{#1}}
\newcommand{\repo}[1]{{\footnotesize\textcolor{gray}{[\,#1\,]}}}
\newcommand{\open}[1]{{\textbf{\textcolor{red!70!black}{[OPEN:\,#1]}}}}
\newcommand{\Zk}{\mathbb{Z}_{2^{k}}}
\newcommand{\Ztfour}{\mathbb{Z}_{2^{64}}}
\newcommand{\thsmc}{\code{threshold\_SMC}}
\newcommand{\initsmc}{\code{init\_SMC}}
\newcommand{\SigmaT}{\Sigma_{T}}

\newtheorem{conjecture}{Conjecture}
\theoremstyle{definition}
\newtheorem{obligation}{Proof obligation}

\lstdefinestyle{mpc}{
  basicstyle=\ttfamily\footnotesize,
  keywordstyle=\color{blue!60!black},
  commentstyle=\color{green!40!black}\itshape,
  showstringspaces=false,
  breaklines=true,
  columns=fullflexible,
  frame=single,
  framesep=4pt,
  xleftmargin=6pt,
}
\lstset{style=mpc}

\title{\textbf{Executable Conformance for a Secure-Multiparty Application:\\
What Implementation-Level Evidence Can Certify, and\\
Where It Stops Short of a Privacy Proof}\\[4pt]
{\normalsize\textsc{First draft --- artifact baseline commit \code{305a3a8}}}}

\author{Harshit Kargeti\thanks{Author list, affiliations, and acknowledgments are
\open{author metadata}. This manuscript was produced through a three-party
adversarial loop (direction, implementation, adversarial review) described in
Section~\ref{sec:process}; the division of credit is not yet fixed.}}
\date{\today}

\begin{document}
\maketitle

\begin{abstract}
Secure-multiparty computation (MPC) \emph{applications}---programs written on top
of an MPC framework---increasingly ship with continuous-integration checks that
are easy to mistake for privacy guarantees. We report a methodology and a
boundary finding for one such application: the deterministic-query fragment of the
knowledge-threshold belief computation of Mardziel, Hicks, Katz and Srivatsa
(PLAS 2012), implemented in MP-SPDZ. Our first contribution is an
\emph{executable-conformance} methodology: an independently written plaintext
oracle, a tiny hand-verifiable fixture containing both an accepted and a rejected
invocation, an anti-echo interface contract that keeps expected outputs out of the
circuit, and SHA-bound, recomputed raw evidence. This makes a green check mean
that a \emph{defined} functionality agrees with an \emph{independent} reference on
the inputs tested, rather than that two co-derived models agree. Our second
contribution is a negative, boundary result. We attempted to \emph{certify}, by
static inspection of the compiled MP-SPDZ assembly, that each recipient's verdict
is delivered privately. Six escalating adversarial re-reviews each produced a
concrete build that leaked a verdict yet passed the then-current gate; the sixth
cannot be closed, because the secret is routed into a comparison subtape through
the very register-argument channel the honest code uses, and because a masked
comparison-open and a raw secret-reveal are the same opcode. We therefore scope
the static check honestly as a \emph{linter} and articulate the hand-off: a real
non-leakage guarantee needs privacy of the executed circuit \emph{and} a semantic
source-to-spec binding, neither of which functional conformance supplies in
general. We are careful not to over-claim: we show the \emph{studied
syntactic/opcode-level approach} is unsound against a source-controlling author,
and state a stronger impossibility only as a conjecture with its missing proof
obligations named. The artifact is public and reproducible against a pinned
backend.
\end{abstract}

% ===========================================================================
\section{Introduction}
\label{sec:intro}

An MPC framework such as MP-SPDZ~\cite{keller2020mpspdz} compiles a high-level
program into a protocol that several parties execute so that each learns only its
prescribed output. Building a correct \emph{application} on top of such a framework
is a separate problem from the framework's own security proof: the application
author can, by accident or by intent, write source that reveals a secret the
functionality was meant to protect. When such an application is wired into
continuous integration, a green check is tempting to read as evidence that the
privacy property holds. This paper is about what that check can and cannot
establish.

Our vehicle is the knowledge-threshold belief computation proposed by Mardziel,
Hicks, Katz and Srivatsa at PLAS 2012~\cite{plas2012}. In their setting, $N$
parties each hold a secret; before releasing a query's output, the parties check
\emph{inside} the MPC whether releasing it would let any party's posterior belief
about another party's secret exceed a threshold, and---crucially---whether a party
receives output or a rejection must not itself be observable by the others. The
authors implemented one of their two mechanisms and \emph{simulated} the other,
``SMC belief tracking,'' without implementing it. We build the smallest faithful
circuit for a deterministic-query fragment of that second mechanism, not as an end
in itself, but as a concrete artifact on which to study conformance and its
limits.

We are explicit about what this paper is \emph{not}. It is not ``the first
implementation of PLAS belief tracking,'' and it makes no such claim: an earlier
round of adversarial review on this project established, correctly, that
implementing and timing a 2012 construction is not by itself a research
contribution absent a new protocol, theorem, compiler, security result, or
empirically supported systems insight.\repo{docs/review-log.md R-10} The
contributions here are methodological and boundary-finding, and the artifact is
their vehicle.

\paragraph{Contributions.}
\begin{enumerate}[leftmargin=1.4em,itemsep=2pt]
\item \textbf{An executable-conformance methodology for an MPC application}
(Section~\ref{sec:method}): an independent oracle, an accept-and-reject fixture, an
anti-echo interface, and recomputed SHA-bound evidence, so that CI checks a defined
functionality against an independent reference.
\item \textbf{A reproducible conformance artifact} for the deterministic-query
fragment of \thsmc{} ($N{=}3$, domain $\{0,1,2\}$, thresholds $1/2$), with 92
conformance tests, a 228-case adversarial coverage sweep, and five committed
executable negative controls, all green against a pinned backend
(Section~\ref{sec:results}).
\item \textbf{An empirical boundary result with a taxonomy}
(Section~\ref{sec:bypass}): six escalating static-delivery bypasses plus an
independent semantic bypass, each reproduced end to end, showing that the studied
opcode-level certification approach is unsound against a source-controlling author,
and isolating the two structural reasons.
\item \textbf{A precise articulation of the hand-off}
(Section~\ref{sec:handoff}): the two separate properties a real non-leakage
guarantee needs, and why conformance supplies neither in general.
\item \textbf{The adversarial review process as a reusable discipline}
(Section~\ref{sec:process}), with an append-only, auditable log including a
self-refutation.
\end{enumerate}

We foreground the honest scope of the boundary result now, so it is not
over-read. What we \emph{demonstrate} is empirical: six concrete builds and a
structural argument. What we do \emph{not} assert is a general impossibility
theorem for static analysis; Section~\ref{sec:conjecture} states the stronger
claim only as a conjecture and lists what a proof would require. A dataflow-aware
analysis is explicitly not conjectured impossible---indeed it is one of the routes
to a real guarantee.

% ===========================================================================
\section{Background}
\label{sec:background}

\paragraph{The functionality.}
Following PLAS 2012~\cite{plas2012}, $N$ parties hold secrets $s_1,\dots,s_N$ in a
finite domain $D$. Each observer $j$ maintains a belief $\delta_j$, a distribution
over the other parties' secrets. A knowledge-threshold policy for party $i$ with
public threshold $t_i$ demands that no other party ever assign posterior
probability greater than $t_i$ to any single value of $s_i$. For a query $Q$, the
mechanism runs, per recipient $j$, a check \code{tcheck} that quantifies over
\emph{every possible output} $o$ of $Q$: it revises $\delta_j$ by conditioning on
$Q(s){=}o$ and rejects if any protected party's marginal would exceed its
threshold. Quantifying over all possible outputs---not merely the actual one---is
what makes the accept/reject decision simulatable and stops the rejection itself
from leaking.\repo{conformance/CONTRACT.md; Fig.~4} Acceptance uses ${\le}\,t_i$
(equality allowed); the boundary is strict~$>$.

We study the fragment in which queries are \emph{deterministic}, total, and public,
and their choice is independent of the secrets. Two queries instantiate it:
\code{p1\_is\_max} ($x_0 \ge x_1 \wedge x_0 \ge x_2$) and \code{sum\_even}
($x_0{+}x_1{+}x_2$ even). Probabilistic queries need likelihood weighting and are
out of scope.\repo{conformance/oracle.py; CONTRACT.md}

\paragraph{Private delivery.}
On accept, recipient $P_j$ learns the actual output $o$; on reject it learns a
rejection, and its belief is unchanged. Each recipient's pair
$(\mathrm{accept}_j, \mathrm{payload}_j)$ must be delivered privately, with the
payload masked to a fixed constant on reject, and no party may learn another
party's verdict.\repo{conformance/INTERFACE.md}

\paragraph{The backend.}
Circuits are compiled and run with MP-SPDZ pinned at commit
\code{9d809599\dots}, using \code{replicated-ring-party.x}: three-party replicated
secret sharing (Rep3) over $\Ztfour$, semi-honest, honest-majority.\repo{.github/workflows/benchmark.yml;
ADVERSARY.md} The MP-SPDZ compiler emits a manifest of \emph{tapes} (a main tape
plus comparison subtapes for equality/less-than), each a sequence of opcodes. Two
facts about this compiled form drive Section~\ref{sec:bypass}: a public reveal and
a legitimate masked comparison-open are the \emph{same} opcode
(\code{asm\_open}/\code{vasm\_open}), and a value computed in one tape reaches
another only through memory or through register arguments passed by
\code{call\_tape}/\code{call\_arg}---a cross-tape register reference is a compiler
error.

\paragraph{Threat model and scope box.}
We assume at most one semi-honest corruption, no collusion, authenticated encrypted
channels, and process isolation. The test harness, however, runs all three parties
as one operating-system user on one host: the isolation assumption is stated, not
enforced by the test.\repo{conformance/ADVERSARY.md} We make no claim under
collusion, a malicious party, unencrypted channels, or a shared-host adversary, and
no performance claim of any kind.

% ===========================================================================
\section{Method: Executable Conformance for an MPC Application}
\label{sec:method}

The discipline has four parts, each of which exists to defeat a specific way a
conformance check can be vacuous.

\paragraph{An independent oracle.}
The plaintext oracle is written from the paper, independently of the circuit and of
the project's earlier exact-rational reference model, and transcribes the
all-possible-outputs \code{tcheck} semantics directly.\repo{conformance/oracle.py}
Independence matters: comparing the circuit against a reference co-derived from it
would only confirm that two models built from the same misunderstanding agree. The
transcription is not assumed correct---it carried a fatal error (conditioning on the
actual output only, rather than all outputs) that adversarial review caught and
that we corrected, adding a discriminating regression test whose two semantics give
different visible results.\repo{docs/review-log.md R-13;
test\_conformance.py::test\_discriminating\_case}

\paragraph{An accept-and-reject fixture.}
The fixture uses three parties, domain $\{0,1,2\}$, thresholds $1/2$, and two
invocations chosen so that one invocation accepts for all recipients and a later
one rejects for one recipient while the others accept, exercising per-recipient
divergence and---the property the earlier circuits violated---belief preservation
on reject.\repo{conformance/CONTRACT.md} A single fixture is a regression vector,
not a proof of general conformance; we say so, and add targeted tests for the
equality boundary, an unrealized output, an impossible event, a non-uniform prior,
and the threshold range.

\paragraph{An anti-echo interface.}
Before any circuit was written, the interface was fixed to forbid an ``echo''
circuit that simply returns what it is told.\repo{conformance/INTERFACE.md;
docs/review-log.md R-17,R-18} The expected outputs and expected post-state never
enter the circuit---neither compiled in nor as runtime inputs; they exist only in
the external harness, which obtains them independently from the oracle. Circuit
inputs are exactly the public parameters $(N, D, Q, t_i)$ and the secret-shared
secrets and beliefs. The harness runs at least one secret state beyond the fixture,
so a constant-output circuit cannot pass by accident. Beliefs travel as dense,
unnormalized non-negative integer weight vectors in a fixed lexicographic order;
two states are equal iff their weight vectors are proportional. The threshold check
avoids division: for public $t_i = a/b$, marginal $M$, and total $Z$, compare
$b\,M \le a\,Z$. Because MP-SPDZ comparison is \emph{signed}, the no-wraparound
requirement is $B < 2^{k-1}$ with $B = \max(a,b)\cdot S \cdot W$, not the naive
$2^k > B$; the naive bound is wrong and we pin a counterexample.\repo{INTERFACE.md;
circuit\_spec.py; R-22} For the fixture, $B = 54 \ll 2^{63}$.

\paragraph{Recomputed, SHA-bound evidence.}
Every run retains a typed JSONL record: each party's raw stdout/stderr, return
code, launch command, a source hash, a delivery signature, provenance, and channel
status.\repo{conformance/validate\_evidence.py; private\_run.py} The validator
enforces an exact field set and types, re-parses each retained transcript, and, in
CI under \code{-{}-recompute}, requires each record's \code{source\_sha256} and
delivery signature to equal values recomputed independently from the checked-out
source and a fresh compile. It also binds each record to a canonical case by
recomputing the input hash and the oracle verdict and requiring a bijection, so a
replayed duplicate, an unbound input, or a forged verdict value is rejected.
Provenance is real: committed local evidence is labeled unbound
(\code{bound:false}); only CI runs carrying \code{GITHUB\_SHA} are bound, and the
job asserts the MP-SPDZ head equals the pin.\repo{R-34,R-37,R-43;
benchmark.yml} A dedicated \code{pipefail} canary proves a failing suite cannot go
green through a logging pipe.\repo{R-31}

\subsection{Artifact and results}
\label{sec:results}
The methodology is instantiated as a public artifact against the pinned backend.
Functional conformance holds on the fixture (two invocations, accept then reject)
and two additional valid states, matching the oracle on verdicts, payloads, and
reconstructed weights (4/4). The adversarial coverage sweep passes 228 of 228
cases (216 single-invocation and 12 carried), spanning all 27 secrets, both
queries, and uniform, non-uniform, and per-party-scaled integer weights near the
bit bound.\repo{coverage.py; results-coverage.txt} The test suites---92
conformance tests (46 functional, 46 private-gate) and 9 reference tests---are
green, as is the CI job that builds MP-SPDZ at the pin, asserts head equals pin,
runs the harness, coverage, and private-delivery gate, and validates the uploaded
evidence under \code{-{}-require-bound -{}-recompute}. The five committed negative
controls of Section~\ref{sec:bypass} are rejected by the linter. No timing is
reported: the repository's older timings measure a different, non-conforming
circuit and are explicitly not to be cited.

% ===========================================================================
\section{The Private-Delivery Gate and Its Adversarial Refutation}
\label{sec:bypass}

Functional conformance says nothing about \emph{privacy} of delivery: a circuit can
compute the right verdicts and then broadcast them. We therefore attempted a
separate, automatic gate that a build must pass before it can be called private,
and subjected it to adversarial re-review. The gate has three layers: a static
inspection of the compiled delivery, a strict runtime transcript check, and the
recomputed evidence binding of Section~\ref{sec:method}. Layers two and three are
sound for what they check. Layer one is the subject of this section.

\subsection{Six escalating bypasses}
\label{sec:six}

The motivating gross case is a build whose \code{reveal\_to(j)} is changed to a
public \code{reveal()} while keeping the per-party print: it compiles to a public
open of all six verdicts, yet a stdout oracle reports success, because a stdout
check cannot distinguish private output from a public open followed by private
printing.\repo{docs/review-log.md R-41} That moved us from a stdout oracle to
inspecting the compiled delivery, and it is committed as the negative control
\code{threshold\_smc\_leaky.mpc}. Six further re-reviews followed, each a concrete
build that leaked a verdict yet passed the gate that existed when it was found.
Table~\ref{tab:taxonomy} is the taxonomy.

\begin{table*}[t]
\centering
\footnotesize
\renewcommand{\arraystretch}{1.25}
\begin{tabularx}{\textwidth}{@{}l l X l l@{}}
\toprule
\textbf{ID} & \textbf{Mutation} & \textbf{Why it passed the then-current gate} &
\textbf{Closed?} & \textbf{Committed control} \\
\midrule
B0 & public \code{reveal()} in main, keep per-party print &
stdout oracle cannot tell private output from public open + private print &
yes & \code{leaky.mpc} \\
B1 & separate \code{@function\_tape} public reveal + print &
inspector read only the main tape &
yes & \code{subleak.mpc} \\
B2 & name-spoofed \code{EQZ(spoof)} subtape: wrong-player delivery + file sink &
masked allowlist keyed on tape \emph{names} &
yes & \code{namespoof.mpc} \\
B3 & \code{reveal(False)} in an EQZ/LTZ-named subtape &
gate treated \code{False} opens as safe (they are still public) &
yes & \code{openfalse.mpc} \\
B4 & \code{reveal(False)} in an \emph{existing} subtape body via scalar
\code{stms}/\code{ldms} &
multiset pins tape names/counts, not bodies &
yes & \code{openfalse.mpc} \\
B5 & same, crossing the boundary with vectorized \code{vstms}/\code{vldms} &
memory regex matched only scalar/gf2n forms &
yes & \code{openfalse\_vec.mpc} \\
B6 & verdict routed via \code{call\_tape}/\code{call\_arg} register arguments &
rules cover only the memory channel; this is the honest channel &
\textbf{no} & --- (unforbiddable) \\
\bottomrule
\end{tabularx}
\caption{The adversarial mutation taxonomy. B0--B5 were each closed by
strengthening the static gate; B6 cannot be closed without breaking the honest
build, which forced the linter-scoping decision. Each closed row has a committed
executable negative control.\repo{docs/limits.md; review-log.md R-41--R-47}}
\label{tab:taxonomy}
\end{table*}

Each closure sharpened the gate: from main-tape-only to whole-manifest inspection
(B1); from a name-based allowlist to content-based rules (B2); from trusting the
open flag to pinning the exact multiset of comparison subtapes (B3); from names and
counts to forbidding the memory channel a verdict would need to enter a subtape
(B4); and from a scalar-only memory pattern to one covering vectorized and
gf2n forms (B5).\repo{review-log.md R-42--R-46} The clean build stores nothing in
memory in its main tape and touches no memory in its subtapes, so forbidding memory
is a sound restriction with no false positives: memory was an attacker-introduced
channel.

\subsection{Why this is structural, not a patch away}
\label{sec:structural}

B6 does not fit the pattern. The verdict is a register in the main tape; a
cross-tape register reference is a compiler error; so the only ways a verdict can
reach a subtape's open are memory---now forbidden---or the register arguments passed
by \code{call\_tape} and received by \code{call\_arg}. But that register channel is
exactly how the honest equality and less-than subtapes receive their operands.
Forbidding it would reject the honest build. So B6 cannot be closed. Two facts
combine to defeat any purely syntactic gate:\repo{docs/limits.md}
\begin{enumerate}[leftmargin=1.4em,itemsep=1pt]
\item \emph{Open-type ambiguity.} A public reveal and a masked comparison-open are
the same opcode; the \code{True}/\code{False} flag controls only a post-open
correctness check, not privacy. Whether an open is safe depends on whether its
operand is masked---a dataflow property present in the assembly but not determined
by opcode identity alone.
\item \emph{An unforbiddable channel.} A secret can be moved into a subtape's open
through the same \code{call\_arg} register channel the honest subtapes use.
\end{enumerate}
Enumerating channels is therefore a losing game against the root ambiguity in~(1).
We scope layer one honestly as a \emph{linter}: it rejects the five committed
negative controls and catches gross or accidental leaks and regressions, but a
PASS is explicitly not a non-leakage certificate. The runtime banners of the tools
say so.\repo{delivery\_inspect.py; private\_run.py}

\subsection{An independent semantic bypass}
\label{sec:semantic}

On a different axis, a build can deliver $(\mathrm{accept}_j,\mathrm{payload}_j)$
privately to exactly the correct recipient and still compute the \emph{wrong}
secret function inside. Privacy of delivery says nothing about \emph{which} value
is delivered. Functional conformance binds source to specification, but only up to
test coverage, not in general. This ``correct recipient, wrong function'' gap is
orthogonal to the six leak channels and is why privacy alone is
insufficient.\repo{docs/limits.md}

\subsection{Can the boundary be made a theorem?}
\label{sec:conjecture}

The directive under which this work proceeds is explicit that we must not overstate.
What Sections~\ref{sec:six}--\ref{sec:structural} establish is an empirical
demonstration plus a structural argument, for \emph{one} backend and the
\emph{studied} opcode-level approach. We do not assert a general impossibility. A
stronger claim would be:

\begin{conjecture}[dataflow-free syntactic certification is unsound]
\label{conj}
Fix the pinned MP-SPDZ Rep3 backend and the functionality $F = \thsmc$. Call a
checker \emph{syntactic/dataflow-free} if it is a computable predicate on the
compiled tape-manifest that is invariant under a transformation class $T$ which
preserves opcode identity and tape structure but may relabel an open's operand
(substituting a verdict-derived register for a masked one) and may move a register
between tapes via the \code{call\_arg} channel. Then there exist source programs
$P_{\mathrm{safe}}$ (delivers $F$ privately) and $P_{\mathrm{leak}}$ (leaks a
verdict) whose compiled manifests are $T$-related; hence any such checker $C$
satisfies $C(P_{\mathrm{safe}}) = C(P_{\mathrm{leak}})$ and cannot both accept
$P_{\mathrm{safe}}$ and reject $P_{\mathrm{leak}}$. No dataflow-free syntactic
checker is sound and complete for non-leakage on this backend.
\end{conjecture}

We state this as a conjecture, not a theorem, because three obligations are unmet:

\begin{obligation}
A formal model of the compiled manifest and a rigorous definition of the
transformation class $T$---i.e., a precise formalization of ``does not do
dataflow'' as invariance under operand relabeling and cross-tape \code{call\_arg}
moves.
\end{obligation}
\begin{obligation}
An explicit $T$-related pair $P_{\mathrm{safe}}, P_{\mathrm{leak}}$ with a proof
that they are $T$-equivalent and that the honest build is a fixed point of the gate.
We have empirical instances (the B3/B4 \code{reveal(False)} pairs), not a proof of
$T$-equivalence.
\end{obligation}
\begin{obligation}
A protocol-level definition of ``leak,'' to prove $P_{\mathrm{leak}}$ actually
leaks and $P_{\mathrm{safe}}$ does not.
\end{obligation}

Obligation~3 is the crux: proving the impossibility rigorously requires the very
simulation-based privacy definition whose absence is the boundary this paper is
about. We therefore keep the asserted result empirical and mark
Conjecture~\ref{conj} \open{theorem: formalize $T$, exhibit a proven $T$-equivalent
pair, and supply a protocol-level leak definition}. Crucially, a \emph{dataflow-aware}
analysis is \emph{not} covered by the conjecture and is not claimed impossible;
Section~\ref{sec:handoff} lists it as a route to a real guarantee.

% ===========================================================================
\section{The Hand-off: What Would Make It a Real Guarantee}
\label{sec:handoff}

The boundary is not that privacy is unknowable; it is that the evidence type
changes. A real non-leakage guarantee needs two \emph{separate} properties, and
functional conformance supplies neither in general.\repo{docs/limits.md}

\begin{description}[leftmargin=1em,itemsep=2pt]
\item[Privacy of the executed circuit.] A protocol-level Rep3 simulation argument
(a human MPC specialist), or a formally-verified comparison primitive plus a
dataflow analysis proving every open's operand is masked and no verdict-derived
register flows into any open. This shows the executed functionality is computed
privately.
\item[Semantic source-to-spec binding.] An argument that the circuit computes the
\emph{intended} belief-tracking function, not merely one that agrees with the
oracle on the tested cases. Without this, privacy still permits the ``correct
recipient, wrong function'' build of Section~\ref{sec:semantic}.
\end{description}

A Rep3 proof alone establishes the first and not the second; conformance testing
gives partial, coverage-bounded evidence toward the second and nothing toward the
first. Table~\ref{tab:handoff} places each property with the actor that must supply
it. The practical recommendation is that an artifact should represent this boundary
in its own outputs---our tools print a linter caveat at runtime and the threat-model
document leads its ``not claimed'' list with sound non-leakage by static
inspection---so that a green check is never mistaken for a privacy proof.

\begin{table}[t]
\centering
\footnotesize
\renewcommand{\arraystretch}{1.25}
\begin{tabularx}{\columnwidth}{@{}l X l@{}}
\toprule
\textbf{Property} & \textbf{Evidence type} & \textbf{Supplied by} \\
\midrule
Functional correctness (tested) & conformance vs.\ independent oracle & this work \\
No gross delivery leak & static linter + strict transcript & this work \\
Evidence integrity & recomputed, SHA-bound provenance & this work \\
Executed-circuit privacy & Rep3 simulation or verified dataflow primitive &
\open{expert / future} \\
Semantic source-to-spec binding & correctness argument beyond test coverage &
\open{future} \\
\bottomrule
\end{tabularx}
\caption{The hand-off. The top three rows are what implementation-level evidence
establishes; the bottom two require a protocol-level argument this work does not
provide.}
\label{tab:handoff}
\end{table}

% ===========================================================================
\section{The Adversarial Review Process}
\label{sec:process}

The results above were produced by a three-party loop: direction, implementation,
and adversarial review, with a shared repository as the only mailbox.\repo{docs/workflow.md}
The reviewer is instructed to refute rather than assess, and to default to
``refuted'' when uncertain. Every critique is logged in an append-only review log,
whether it survived or not; the log has 47 entries.\repo{docs/review-log.md} We
report the process as a contribution for two reasons. First, it has teeth: it
caught a fatal specification bug (Section~\ref{sec:method}) and produced the entire
bypass sequence of Section~\ref{sec:bypass}. Second, it is self-correcting in a way
worth documenting---our own soundness premise that ``memory is the only cross-tape
channel'' was asserted at one round (justifying the memory-channel rules) and
refuted at the next by the \code{call\_arg} channel, and both the assertion and its
refutation remain in the log.\repo{review-log.md R-45,R-47} Two models agreeing is
weak evidence, for correlated reasons, so the process is a filter, not a
substitute: the project's standing rule is that a human MPC specialist must sign
off before any security claim ships.

% ===========================================================================
\section{Related Work}
\label{sec:related}

We position against five categories. We flag at the outset that the
novelty statements below are \open{literature review: the project has not run a
search reaching Google Scholar or Semantic Scholar}, and that the paper's framing
depends on completing it.

\paragraph{MPC frameworks and compilers.} MP-SPDZ~\cite{keller2020mpspdz} and peers
provide the substrate; to our current knowledge they do not offer application-level
conformance against an independent oracle nor certify per-recipient delivery
\open{confirm}.

\paragraph{Knowledge-based security and quantitative information flow.} The
functionality is from PLAS 2012~\cite{plas2012}, which simulated rather than
implemented SMC belief tracking; we target its deterministic-query fragment
\open{confirm no later work implements it with a conformance target}.

\paragraph{Formal verification of MPC and secure compilation.} Verified MPC
compilers, type systems for information flow, and computational-soundness results
are the correct heavy machinery our boundary hands off to; the negative result is
positioned as ``why an opcode-level syntactic shortcut fails,'' not as competing
with dataflow-aware verification \open{cite the specific lines of work}.

\paragraph{Conformance, differential, and metamorphic testing of cryptographic
code.} Test-vector conformance and differential testing of crypto libraries are the
closest methodological neighbors; the transfer to an MPC \emph{application} with an
anti-echo interface and recomputed evidence is the claimed novelty
\open{substantiate against that literature}.

\paragraph{Static analysis and linting for information flow.} Our negative result
names a specific failure mode for opcode-level linters of MPC bytecode: opcode
identity versus dataflow, and an honest-used register channel \open{position
against existing IFC-linting work}.

% ===========================================================================
\section{Threats to Validity}
\label{sec:threats}

\emph{Construct.} ``Conformance'' means agreement with an independent oracle on
tested inputs; the oracle is a human transcription of one paper's
deterministic-query fragment and carried a fatal bug until review caught it, so
unreviewed transcription errors may remain, and no MPC expert has yet signed off.
\emph{Internal.} The linter's soundness for the memory channel rests on empirical
facts about the pinned backend (the clean build uses no memory; a cross-tape
register reference is a compile error); a different backend version could change
opcode semantics or channels. The \code{tls\_certs\_present} flag is a weak proxy
(files exist $\ne$ wire encrypted), a caveat we state.\repo{ADVERSARY.md}
\emph{External.} Results are for one functionality, $N{=}3$, domain $\{0,1,2\}$, one
backend, semi-honest honest-majority, and a single-host harness that is not an
adversarially isolated deployment; generalization is \open{E4}. \emph{The boundary
result is empirical, not a theorem} (Section~\ref{sec:conjecture}). \emph{Semantic
gap.} Even perfect delivery privacy would not establish that the circuit computes
the intended function beyond tested cases. \emph{No performance claims} are made;
the repository's timings measure a different, non-conforming circuit and are
explicitly not to be cited.\repo{README; docs/gap.md}

% ===========================================================================
\section{Limitations and Future Work}
\label{sec:future}

Persistent secret-shared state across invocations ($\SigmaT$) is unimplemented and,
under the current directive, unauthorized; it is future work, not a gap we fill
here. The single highest-value next step is the literature review that gates the
novelty framing (Section~\ref{sec:related}). The experiment most likely to change a
conclusion is a genuine \emph{dataflow-aware} delivery checker: if it closes the gap
the linter cannot, it changes the boundary; if it fails, it upgrades the negative
result from ``syntactic gates fail'' to ``here is exactly where dataflow is
needed.'' \open{dataflow checker: new implementation, requires authorization}
Attempting the impossibility theorem of Section~\ref{sec:conjecture}, generalizing
beyond the fixture to larger domains and probabilistic queries, and obtaining the
protocol-level Rep3 argument from a specialist are the remaining directions.

% ===========================================================================
\section{Conclusion}
\label{sec:conclusion}

Executable conformance against an independent oracle, backed by an anti-echo
interface and recomputed evidence, is a rigorous and reusable discipline for MPC
applications: it establishes real functional correctness on tested inputs, catches
gross information leaks, and makes a green check mean something. It also has a
sharp, nameable boundary. Attempting to certify per-recipient non-leakage by static
inspection of the compiled assembly fails against an author who controls the
source, for structural reasons---a masked open and a raw reveal are the same
opcode, and the secret can travel the same register channel the honest code
uses---and that boundary is exactly where implementation-level evidence must hand
off to a protocol-level privacy proof. Naming the boundary, rather than papering
over it with a green check, is the contribution we most want to stand.

% ===========================================================================
\appendix
\section{Claim--evidence--limitation table}
\label{app:claims}

Table~\ref{tab:claims} lists the claims the paper makes, each with its repository
evidence and the limitation stated alongside it.

\begin{table*}[t]
\centering
\footnotesize
\renewcommand{\arraystretch}{1.2}
\begin{tabularx}{\textwidth}{@{}l X X@{}}
\toprule
\textbf{Claim} & \textbf{Evidence} & \textbf{Limitation} \\
\midrule
C1 Circuit reproduces the oracle on an accept-and-reject fixture &
\code{harness.py}, \code{oracle.py}, 4/4 named cases &
per-tested-input, not general; one fixture is a regression vector \\
C2 Agreement across a 228-case adversarial sweep &
\code{coverage.py}; \code{216 single + 12 carried = 228 pass} &
fixed enumerable set at $N{=}3$, $D{=}\{0,1,2\}$ \\
C3 Oracle is independent and transcribes all-outputs \code{tcheck} &
\code{oracle.py}, R-13 fix, discriminating test &
human transcription; no expert sign-off \\
C4 Interface forbids echo-cheating &
\code{INTERFACE.md}, extra secret state &
enforced by construction/review, not a machine proof \\
C5 Signed bound is $B<2^{k-1}$; fixture safe &
\code{circuit\_spec.py}, R-22 &
general bound stated, not auto-checked \\
C6 Linter rejects five negative controls &
\code{delivery\_inspect.py}, five \code{.mpc} controls &
a linter, not a non-leakage proof \\
C7 Runtime transcript fails closed &
\code{strict\_parse\_party}, \code{test\_private.py} &
cannot see shares/timing/network \\
C8 Evidence is typed, SHA-bound, recomputed &
\code{validate\_evidence.py}, CI \code{-{}-recompute} &
does not bind source to intended spec beyond tested cases \\
C9 Provenance is real (bound only in CI) &
\code{repo\_provenance()}, R-34 &
not a reproducibility guarantee on arbitrary machines \\
C10 Studied syntactic gate is unsound vs.\ a source-controlling author &
R-41--R-47, \code{docs/limits.md} &
empirical + structural, not a theorem; this backend \\
C11 Conformance binds source to spec only up to coverage &
\code{docs/limits.md}, \code{gap.md} &
motivates the semantic hand-off property \\
C12 The review loop is auditable and self-correcting &
\code{review-log.md} R-45$\to$R-47 &
process evidence, not completeness \\
\bottomrule
\end{tabularx}
\caption{Claim, evidence, and limitation---grounded only in the repository at
\code{305a3a8}.}
\label{tab:claims}
\end{table*}

\section{A minimal leak and the honest channel}
\label{app:listings}

The \code{reveal(False)} bypass (B3/B4) reconstructs a verdict publicly with the
``correctness-check-disabled'' flag, which the earlier gate wrongly treated as
safe:

\begin{lstlisting}[language=Python]
# spoofed EQZ(3)_63 subtape: no stdout, no privateoutput, no sink
w = sint.load_mem(4000)   # verdict moved from main
w.reveal(False)           # STILL a public open; False != private
\end{lstlisting}

\noindent B6 is different: the honest comparison subtapes receive their operands
through the same register channel a leak would use, so it cannot be forbidden:

\begin{lstlisting}[language=Python]
# main tape invokes a comparison subtape by passing registers directly:
#   call_tape 12, 3, ci0, ... s6(3), s0(3) ...   (no memory)
# the subtape receives them via call_arg -- the honest path.
\end{lstlisting}

\begin{thebibliography}{9}
\bibitem{plas2012}
P.~Mardziel, M.~Hicks, J.~Katz, and M.~Srivatsa.
\newblock Knowledge-Oriented Secure Multiparty Computation.
\newblock In \emph{PLAS}, 2012.
\newblock \open{full citation and DOI/URL to be finalized}.

\bibitem{keller2020mpspdz}
M.~Keller.
\newblock MP-SPDZ: A Versatile Framework for Multi-Party Computation.
\newblock In \emph{ACM CCS}, 2020.
\newblock \open{confirm exact citation details}.
\end{thebibliography}

\end{document}
